\section{The actual third-growth pair in the source clock and frame}
\label{tgframe:section}

This calculation uses equations (6.2)--(6.12), (7.2)--(7.12), and
Lemma 7.4 of the retained 166-page source. Its SHA-256 is displayed
in two consecutive segments:
\[
\begin{gathered}
\texttt{0e779481c4da40bd28d1e642e1d8ca57}\\
\texttt{447d129610df28dfa5a11e9af8ae228f}.
\end{gathered}
\]
The coupled input is the actual third-growth interval proved in
\path{finite_stage/third_growth_entry.tex}. The source's pulse
equation is written at physical viscosity one. We retain an arbitrary
positive physical viscosity $\nu_{\rm NS}$ when computing a candidate
residual; when identifying the source's particular homogeneous pulse,
we explicitly set $\nu_{\rm NS}=1$ and harmonic $m=1$.
The original coupled time scale $\nu=\sqrt{A_0}$, the physical coupled
diffusivity $d$, and the source chart parameter $\varepsilon=Q^h$
are distinct.
The source growth constant denoted $\lambda_0$ in $\lambda(v)$ below
is the source's own growth constant. It is not identified with the
coupled integer frequency $\lambda_0=\mu=\lceil\sqrt{A_0}/2\rceil$.

\subsection{Exact physical differentiation of the source pulse clock}

Fix a source label and one of its local lifts. Set
\[
 b_g=\sqrt2-1,\quad v_r=(1,-b_g),\quad v_t=(b_g,1),\quad
 J_g=\begin{pmatrix}3&1\\1&5\end{pmatrix},\quad
 T_g=4+\sqrt2,\quad\Lambda_g=4-\sqrt2.
\]
Direct multiplication gives $J_gv_t=T_gv_t$ and
$J_gv_r=\Lambda_gv_r$, while $v_t\cdot v_r=0$ and
$|v_t|^2=|v_r|^2=1+b_g^2$. The source retains
\[
\begin{gathered}
 d_r=2\left((1+h)\frac{\log\Lambda_g}{\log T_g}-h\kappa_s\right),\qquad
 \kappa_s=10^{-5},\\
 Q=2^{-\ell},\quad
 i=\left\lfloor\log_{T_g}\frac{Q^{-1-h}}{S_*}\right\rfloor,\quad
 c_i=T_g^iQ^{1+h}.
\end{gathered}
\]
On the local lift
\[
 Y_i-c_\gamma-k_{\rm copy}=\xi_gv_r+\eta_gv_t,\qquad
 v=\frac{\eta_g+r_0}{c_i},\qquad L_s=\frac{2r_0}{c_i},
 \quad
 \eta_g=\frac{v_t\cdot(Y_i-c_\gamma-k_{\rm copy})}{1+b_g^2}.
\]
The label and integer $k_{\rm copy}$ are held fixed on this lift.
Consequently $L_i\eta_g=0$, $N_i\eta_g=1$, where
$L_i=v_r\cdot\partial_{Y_i}$ and $N_i=v_t\cdot\partial_{Y_i}$.
The clock has no explicit dependence on the slow variables $R,Z,T$.
The evaluated chart derivatives are exactly
\[
 D_r=\partial_R+M_id_rR^{d_r-1}L_i,\qquad
 D_z=\varepsilon\partial_Z,\qquad
 t_*=-\varepsilon\partial_T+c_iN_i,
 \quad M_i=\Lambda_g^iQ^{d_r/2}.
\]
It follows by applying these displayed operators, without estimating
any coefficient, that
\begin{equation}\label{tgframe:clockchart}
 D_rv=D_zv=0,\qquad t_*v=1.
\end{equation}
There is also an explicit physical verification. Evaluate at the
source map $Y=v_rr^{d_r}+v_tt$ modulo the lattice. Locally the lift of
$Y_i=J_g^iY$ differs from
$\Lambda_g^iv_rr^{d_r}+T_g^iv_tt$ by a constant lattice vector.
The dual formula for $\eta_g$ therefore gives
$\eta_g=T_g^it+C_\gamma$ with a locally constant $C_\gamma$.
Thus, on the evaluated local rectangle,
\begin{equation}\label{tgframe:clockphysical}
 \partial_rv=\partial_zv=\partial_\theta v=0,\qquad
 \partial_tv=\frac{T_g^i}{c_i}=Q^{-1-h}.
\end{equation}
The $\partial_t$ in this equation is the derivative of the evaluated
physical function. On the extended domain the corresponding derivative
is $\partial_t+N_{\rm abs}$. These statements are identical after
evaluation; neither drops the auxiliary derivative.

Let $t_b$ and $r_a>0$ be the actual second return time and the actual
third target-crossing elapsed time. Define
\begin{equation}\label{tgframe:coupledclock}
 t_c(v)=t_b+\alpha v,\qquad \alpha=\frac{r_a}{L_s}>0.
\end{equation}
This parametrizes the entire proved coupled growth interval, with its
original endpoints. In physical source variables its derivative is
$\partial_tt_c=\alpha Q^{-1-h}$ and all three spatial derivatives vanish.
The positive orientation of this clock asserts no sign for a velocity.
The coupled amplitudes below are strictly negative.

\subsection{The actual frame, its inverse, and the direct candidate}

Write the original source phase normal as
\[
 n=\left(x_0-v(pF_R+p_zG_R),\ \frac pR,\
             p_z-\varepsilon v(pF_Z+p_zG_Z)\right).
\]
All $p,p_z,x_0,k,Q$ are the original fixed label constants, including
$kp\in\mathbb Z\setminus\{0\}$. Set
\[
\begin{gathered}
 K_a=\frac{n_{\tan}}{|n_{\tan}|},\quad
 N_a=((K_a)_z,-(K_a)_\theta),\quad
 s_a=\frac{n_r}{|n_{\tan}|},\\
 c(v)=c_0\sqrt{1+s(v)^2},
 \quad
 s(v)=\sigma\left(\frac{u_*}{2}+\frac{u_*v}{L_s}\right).
\end{gathered}
\]
Here $c_0<0$ is the unchanged source constant, and $K_a,N_a$
are identified with their tangential three-vectors. In particular
$c(v)\ne0$. Because $p/R\ne0$, every formula is defined throughout
the source rectangle. The actual source frame is
\[
 U=[e_r-s_aK_a,\ N_a],\qquad
 M=\begin{pmatrix}1&1\\c&-c\end{pmatrix},\qquad B=UM.
\]
The two columns of $U$ are orthogonal, with squared lengths
$g_a^2=1+s_a^2$ and $1$. Both are orthogonal to $n$.
Since $\det M=-2c\ne0$, $B$ is a bijection from $\mathbb R^2$
onto $n^\perp$. The source left inverse, defined on all of
$\mathbb R^3$, is
\begin{equation}\label{tgframe:leftinverse}
 B_\ell=M^{-1}\begin{pmatrix}e_r^T\\N_a^T\end{pmatrix},
 \qquad
 M^{-1}=\frac12\begin{pmatrix}1&c^{-1}\\1&-c^{-1}\end{pmatrix}.
\end{equation}
Indeed the last two rows applied to $U$ give $I_2$.

The eigenvalues of
\[
 B^TB=
 \begin{pmatrix}g_a^2+c^2&g_a^2-c^2\\
                 g_a^2-c^2&g_a^2+c^2\end{pmatrix}
\]
are $2g_a^2$ and $2c^2$, with orthogonal eigenvectors $(1,1)$
and $(1,-1)$. The Euclidean operator norm and the inverse norm on
the tangent plane are therefore exactly
\begin{equation}\label{tgframe:framenorms}
 \|B\|=\sqrt2\max(g_a,|c|),\qquad
 \|(B|_{\mathbb R^2})^{-1}\|_{n^\perp\to\mathbb R^2}
       =\frac1{\sqrt2\min(g_a,|c|)}.
\end{equation}
The ambient norm of the particular extension \eqref{tgframe:leftinverse}
is $(1/\sqrt2)\max(1,|c|^{-1})$; it need not equal the inverse norm
restricted to $n^\perp$.

Retain the actual coupled data
\[
 K_3=\mu\widehat\lambda_3,\quad
 S_3=\frac{A_0}{\mu}\widehat\lambda_3^{-k_3-6},\quad
 P_3^*=\frac{A_0}{\mu}\widehat\lambda_3^{-7/8}=S_3e^{L_3},
 \quad u_i=\sqrt{\frac{b_i}{a_i}}
           =\frac{K_3}{\sigma_2\sqrt{n_i}}.
\]
On the interval \eqref{tgframe:coupledclock}, the already constructed
solution is
\[
 y_3(t_c)=\binom{\Theta_3}{\Omega_3}
        =-P(v)\binom{1}{u(v)},\qquad
 S_3\le P(v)\le P_3^*,\qquad
 \frac{u_i}{\sqrt2}\le u(v)\le2u_i,
\]
with $P(0)=S_3$, $P(L_s)=P_3^*$ and $u(0)=u_i$.
These inequalities follow from the actual evolving-parent system,
rather than from a new frame approximation.
The direct candidate $t_{\rm dir}=By_3$ consequently has the exact
components
\begin{align}
 (t_{\rm dir})_r&=-P(1+u),\nonumber\\
 (t_{\rm dir})_{\tan}
       &=P s_a(1+u)K_a-Pc(1-u)N_a,\label{tgframe:directcomponents}\\
 |t_{\rm dir}|^2
       &=P^2\bigl[(1+s_a^2)(1+u)^2+c^2(1-u)^2\bigr],\nonumber\\
 (t_{\rm dir})_r(t_{\rm dir})_{\tan}
       &=-P^2s_a(1+u)^2K_a+P^2c(1-u^2)N_a.\nonumber
\end{align}
These follow by multiplying $My_3=-P(1+u,c(1-u))^T$ and using
the orthogonality of the two columns of $U$. In particular
\[
 |(t_{\rm dir})_r|\ge S_3(1+u_i/\sqrt2),\qquad
 |t_{\rm dir}|\le\sqrt2\max(g_a,|c|)P_3^*\sqrt{1+4u_i^2}.
\]
At the actual target,
\begin{equation}\label{tgframe:endpointomega}
 \frac{P_3^*u_i}{\sqrt2}\le|\Omega_3(t_b+r_a)|
                      \le2P_3^*u_i,\qquad
 P_3^*u_i=\frac{A_0}{\sigma_2\sqrt{n_i}}
                          \widehat\lambda_3^{1/8}.
\end{equation}
The last equality retains $\mu$ until it cancels against its identical
factor in $K_3$. Thus a decreasing temperature target does not give a
decreasing two-component target.

\subsection{The complete principal mismatch, including the cutoff}

With a prime denoting $v$ differentiation at fixed slow variables,
retain
\[
 \mathsf K=\begin{pmatrix}0&-2F&0\\2F+RF_R&0&0\\G_R&0&0\end{pmatrix},
 \quad
 A_\Phi=-\mathsf K+
        \frac{n(n^T\mathsf K-(n')^T)}{|n|^2},
 \quad H=B_\ell(A_\Phi B-B').
\]
Differentiating $n^TB=0$ shows
$n^TB'=-(n')^TB$. Direct multiplication gives
$n^TA_\Phi B=-(n')^TB$ as well. Therefore
$A_\Phi B-B'$ takes values in $n^\perp$, and
\begin{equation}\label{tgframe:frametransport}
 A_\Phi B-B'=BH.
\end{equation}
The source defines the exact error matrix
$E=H-\operatorname{diag}(\lambda,-\lambda)$ with
$\lambda=\lambda_0/\sqrt{1+s^2}$. Its source estimate
$\|E\|\le C/S_*$ is not an assertion that $E$ is zero.

The actual coupled matrix and arbitrary-viscosity source loss are
\[
\begin{gathered}
 C_c=\begin{pmatrix}-d_c&a_3\\b_3&-d_c\end{pmatrix},
 \quad d_c=dK_3^2R_3,\quad
 a_3=\frac{N_3}{K_3R_3},\quad b_3=\frac{K_3c_{\rm det}}{\chi_{\rm old}},\\
 \quad
 \delta_\nu=\nu_{\rm NS}m^2\varepsilon k^2|n|^2,\quad
 \Delta_\nu=\delta_\nu-\alpha d_c.
\end{gathered}
\]
Here $c_{\rm det}$ and $\chi_{\rm old}$ are the third-growth determinant
constant and older phase length; these names distinguish them from the
source $c(v)$ and the cutoff below. Every coupled quantity is evaluated
at $t_c(v)$. The actual third growth has zero added amplitude controls,
so $y_3'=\alpha C_cy_3$ exactly.

Let $\chi=\eta(R,Z,T)\chi_g(\xi_g)\psi(v)$ be a smooth scalar
cutoff with the source support, and set $t_\chi=\chi By_3$.
The homogeneous-pressure formula, with zero principal source, is
\[
 \pi_\chi=\frac{i}{km|n|^2}
        \bigl(n\cdot\mathsf Kt_\chi-n'\cdot t_\chi\bigr).
\]
The cutoff leaves $n^Tt_\chi=0$. Differentiating this identity
and substituting the pressure gives
\begin{equation}\label{tgframe:residual}
 t_\chi'+\mathsf Kt_\chi+\delta_\nu t_\chi+ikmn\pi_\chi
 =B\mathcal E_\chi,\qquad
 \mathcal E_\chi=\chi(\alpha C_c-H+\delta_\nu I)y_3+\chi'y_3.
\end{equation}
For completeness, the normal component of the sum before pressure is
$-n'\cdot t_\chi+n\cdot\mathsf Kt_\chi$; multiplication by $ikmn$
times the displayed pressure cancels it. The remaining vector is
$t_\chi'-A_\Phi t_\chi+\delta_\nu t_\chi$. Use
\eqref{tgframe:frametransport} and $y_3'=\alpha C_cy_3$ to obtain
\eqref{tgframe:residual}.

In original scalar entries this is
\begin{equation}\label{tgframe:residualentries}
 \mathcal E_\chi=-P
 \left\{\chi
 \binom{\Delta_\nu-\lambda-E_{11}+(\alpha a_3-E_{12})u}
       {\alpha b_3-E_{21}+(\Delta_\nu+\lambda-E_{22})u}
       +\chi'\binom1u\right\}.
\end{equation}
All entries are signed. In particular the damping mismatch is
$\delta_\nu-\alpha dK_3^2R_3$, not the sum of the losses.
At fixed slow variables and fixed $\xi_g$, $\chi'=\eta\chi_g\psi'$.
The physical derivative of $\chi$ also includes
$-\varepsilon(\partial_T\eta)\chi_g\psi$ in $t_*\chi$; that slow term
belongs to the complete physical residual and has not been included in,
or removed by, the principal equation \eqref{tgframe:residual}.
A direct finite bound is
\[
 |\mathcal E_\chi|
 \le P\sqrt{1+u^2}\left[
 |\chi|\bigl(|\Delta_\nu|+\lambda+\|E\|
                  +\alpha\max(a_3,b_3)\bigr)+|\chi'|\right].
\]
This is the triangle inequality and the exact norm
$\left\|\begin{smallmatrix}0&a_3\\b_3&0\end{smallmatrix}\right\|
=\max(a_3,b_3)$, with $a_3,b_3>0$.
It retains the carrier, frame, cutoff and all physical diffusion factors.

\subsection{An explicit map of the actual seed to the actual source pulse}

For this subsection only, take $m=1$, $\nu_{\rm NS}=1$, so
$\delta_1=\varepsilon k^2|n|^2$. This is precisely the source
homogeneous pulse equation with its given background. Put
\[
 d_{\rm ref}=\varepsilon k^2B_s^2(1+s^2),\qquad
 {\cal P}_\sigma(v)=
   \exp\left(\int_{L_s/2}^{v}(\lambda(w)-d_{\rm ref}(w))\,dw\right).
\]
Every coefficient is finite on the closed interval. Thus
$\mathcal P_\sigma(0)>0$ exactly, without a Gaussian approximation.
Source Lemma 7.4 specifies the initial coordinate vector
\[
 z_{\sigma,0}=\binom{{\cal P}_\sigma(0)}0,\qquad
 t^h_\sigma(0)=B(0)z_{\sigma,0}.
\]
The fixed source label may affect this strictly positive scalar.

For any nonzero real target vector $z_0=(z_1,z_2)^T$, define
\begin{equation}\label{tgframe:L0general}
 L_0=-\frac1{S_3(1+u_i^2)}
 \begin{pmatrix}z_1&-z_2\\z_2&z_1\end{pmatrix}
 \begin{pmatrix}1&u_i\\-u_i&1\end{pmatrix}.
\end{equation}
This is an explicit invertible linear map with
\[
 L_0[-S_3(1,u_i)^T]=z_0,\qquad
 \det L_0=\frac{|z_0|^2}{S_3^2(1+u_i^2)}>0,
\]
and inverse
\begin{equation}\label{tgframe:L0inverse}
 L_0^{-1}=-\frac{S_3}{|z_0|^2}
       \begin{pmatrix}1&-u_i\\u_i&1\end{pmatrix}
       \begin{pmatrix}z_1&z_2\\-z_2&z_1\end{pmatrix}.
\end{equation}
Indeed the second matrix of \eqref{tgframe:L0general} sends
$(1,u_i)^T$ to $(1+u_i^2,0)^T$, and each matrix times its
transpose is a positive scalar times $I_2$.
The same multiplication verifies the inverse and proves the exact norms
\begin{equation}\label{tgframe:L0norms}
 \|L_0\|=\frac{|z_0|}{S_3\sqrt{1+u_i^2}},\qquad
 \|L_0^{-1}\|=\frac{S_3\sqrt{1+u_i^2}}{|z_0|}.
\end{equation}
In particular the source target gives
\[
 L_0=-\frac{{\cal P}_\sigma(0)}{S_3(1+u_i^2)}
        \begin{pmatrix}1&u_i\\-u_i&1\end{pmatrix},
 \quad
 \|L_0\|=
 \frac{{\cal P}_\sigma(0)\,\mu\widehat\lambda_3^{k_3+6}}
 {A_0\sqrt{1+K_3^2/(\sigma_2^2n_i)}}.
\]
No seed or source initial amplitude has been set to one.

Define the two identity-initial-value fundamental matrices by
\[
 X_c'=\alpha C_cX_c,\quad X_c(0)=I_2,\qquad
 X_s'=(H-\delta_1I_2)X_s,\quad X_s(0)=I_2.
\]
Here is a direct existence, invertibility and uniqueness construction.
For a continuous matrix $D$ on $[0,L_s]$ let
\[
 X_D(v)=I_2+
 \sum_{j\ge1}\int_{0\le w_j\le\cdots\le w_1\le v}
       D(w_1)\cdots D(w_j)\,dw_j\cdots dw_1 .
\]
If $\sup\|D\|=M_D$, the $j$th term is bounded by
$(M_DL_s)^j/j!$. Uniform convergence and regrouping the integrals
give $X_D=I_2+\int_0^vD(w)X_D(w)\,dw$, hence the differential
equation. The corresponding right-ordered series for
$Z'=-ZD$, $Z(0)=I_2$, converges by the same bound.
Differentiation gives $(ZX_D)'=0$, hence $ZX_D=I_2$.
For square matrices this is invertibility. If two solutions agree
initially, the difference satisfies the zero-initial-value integral
equation; iteration bounds its norm by a fixed finite bound times
$(M_Dv)^j/j!$ for every $j$, which tends to zero.
This proves uniqueness. Smooth coefficient dependence follows from
the same series with the following explicit bound. On a compact
coefficient-parameter neighborhood let $M_N\ge1$ bound every coefficient
derivative through order $N$. Differentiating a product of $j$ factors
$N$ times allocates each derivative to one of $j$ factors, so there are
at most $j^N$ terms, each of norm at most $M_N^j$. The differentiated
$j$th integral is therefore bounded by
$j^N(M_NL_s)^j/j!$, a summable sequence. Uniform convergence of these
derivative series justifies termwise differentiation: integrate each
first-derivative series along a parameter segment and interchange the
uniform sum with that integral, then iterate through order $N$.
The right-ordered inverse series has the same bound. Pulse derivatives
then follow from the differential equations. This proves smoothness,
with the explicit derivative equations below providing further bounds.

Now set
\begin{equation}\label{tgframe:intertwiner}
 L(v)=X_s(v)L_0X_c(v)^{-1},\qquad
 L^{-1}=X_cL_0^{-1}X_s^{-1}.
\end{equation}
Differentiation gives
\[
 L'=(H-\delta_1I_2)L-\alpha LC_c.
\]
The original solution is $y_3(v)=X_c(v)y_3(0)$.
Consequently $L(v)y_3(v)=X_s(v)z_{\sigma,0}$, and
\begin{equation}\label{tgframe:exactpulsematch}
 B(v)L(v)y_3(v)=t^h_\sigma(v)
\end{equation}
by uniqueness of the source homogeneous equation and the exact initial
seed calculation. Thus this finite map reaches the actual source pulse.
It does not merely send the coupled pair into an unspecified transverse
solution. Multiplying both sides by the identical source cutoffs
preserves equality of the amplitudes, and applying identical linear
curl and pressure constructions preserves equality of the resulting
fields. Multiplication by the source prescribed scalar leading-wave
amplitude also preserves that equality. The cutoff principal residual
is then $\chi'BLy_3$; the remaining source construction uses its own
estimates for that exact source object.

\subsection{All scale costs of the finite intertwiner}

Here and below $\|\cdot\|$ is the Euclidean matrix norm.
Put $D_i=\operatorname{diag}(1,u_i)$,
$\kappa_i=\max(u_i,u_i^{-1})$, and
\[
 \beta_c(v)=\frac12\left(a_3(t_c(v))u_i+
                                  \frac{b_3(t_c(v))}{u_i}\right).
\]
The actual coefficient bounds $a_3\le a_i$, $b_3\le b_i$ and
$a_iu_i=b_i/u_i=\gamma_i$ give $0<\beta_c\le\gamma_i$.
Conjugating the coupled matrix by the constant $D_i$ gives
\[
 D_i^{-1}C_cD_i=
 \begin{pmatrix}-d_c&a_3u_i\\b_3/u_i&-d_c\end{pmatrix}.
\]
Its symmetric part has eigenvalues $-d_c\pm\beta_c$.
For any solution $w$ in these coordinates, differentiating $|w|^2$
therefore gives
\[
 2\alpha(-d_c-\beta_c)|w|^2
 \le (|w|^2)'
 \le2\alpha(-d_c+\beta_c)|w|^2.
\]
Multiplication by the corresponding scalar integrating factors proves
the following bounds, including the lower singular-value estimate
needed for the inverse:
\[
 \|X_c(v)\|\le\kappa_i
       e^{\int_0^v\alpha(\beta_c-d_c)\,dw},\qquad
 \|X_c(v)^{-1}\|\le\kappa_i
       e^{\int_0^v\alpha(\beta_c+d_c)\,dw}.
\]
For $D_s=H-\delta_1I_2=
\operatorname{diag}(\lambda,-\lambda)-\delta_1I_2+E$,
its symmetric part is bounded between
$(-\lambda-\delta_1-\|E\|)I_2$ and
$(\lambda-\delta_1+\|E\|)I_2$. The identical energy argument gives
\[
 \|X_s(v)\|\le
 e^{\int_0^v(\lambda-\delta_1+\|E\|)\,dw},\qquad
 \|X_s(v)^{-1}\|\le
 e^{\int_0^v(\lambda+\delta_1+\|E\|)\,dw}.
\]
Combining these four proved estimates with \eqref{tgframe:intertwiner}
yields
\begin{align}
 \|L(v)\|&\le\kappa_i\|L_0\|
 e^{\int_0^v[\lambda-\delta_1+\|E\|
                       +\alpha(\beta_c+d_c)]\,dw},\nonumber\\
 \|L(v)^{-1}\|&\le\kappa_i\|L_0^{-1}\|
 e^{\int_0^v[\lambda+\delta_1+\|E\|
                       +\alpha(\beta_c-d_c)]\,dw}.
 \label{tgframe:intertwinernorms}
\end{align}
All quantities on the right are the retained, finite coefficient paths.
In particular the exact coupled diffusion integral at the endpoint is
\[
 \int_0^{L_s}\alpha d_c\,dv
 =dK_3^2\int_0^{r_a}R_3(t_b+r)\,dr
 \in[dK_3^2r_a,\ 2dK_3^2r_a].
\]
The equality is the substitution $r=\alpha v$; the interval is the
proved actual phase-length bound. The determinant is also exact:
\begin{equation}\label{tgframe:detL}
 \det L(v)=\frac{|z_{\sigma,0}|^2}{S_3^2(1+u_i^2)}
 \exp\left(\int_0^v[
        \operatorname{tr}E-2\delta_1+2\alpha d_c]\,dw\right).
\end{equation}
For a differentiable $2$ by $2$ matrix, expanding its two-column
determinant and inserting $X'=DX$ gives
$(\det X)'=(\operatorname{tr}D)\det X$. Its initial determinant
is one. Apply this calculation to $X_s$ and $X_c$, and use
$\operatorname{tr}H=\operatorname{tr}E$ and
$\operatorname{tr}C_c=-2d_c$ to prove \eqref{tgframe:detL}.

Higher pulse derivatives have the explicit recurrence
\[
 L^{(n+1)}=
 \sum_{j=0}^n\binom nj\left[
 D_s^{(j)}L^{(n-j)}-L^{(n-j)}D_c^{(j)}\right],
 \qquad D_c=\alpha C_c .
\]
For a coefficient parameter derivative $\partial_\xi$ the exact
inhomogeneous equation is
\[
 (\partial_\xi L)'=D_s\partial_\xi L
       -(\partial_\xi L)D_c
       +(\partial_\xi D_s)L-L(\partial_\xi D_c).
\]
Its initial value is the derivative of the explicit
\eqref{tgframe:L0general}. To see directly that this produces a finite
bound, multiply on the left by $X_s^{-1}$ and on the right by $X_c$,
differentiate, and integrate the last two terms with these same
multipliers. The preceding norm estimates bound every resulting
integral on $[0,L_s]$. Iteration proves the analogous assertion for
every finite mixed derivative; all product-rule terms and initial
derivatives are retained. This also proves smoothness on each compact
source slow neighborhood where the frame is defined.

The factor $\mathcal P_\sigma(0)/S_3$ and the two explicit propagation
costs in \eqref{tgframe:intertwinernorms} can be very large or very small.
There is no supplied relation identifying $L_s$ with $L_3$.
Thus these equations prove a finite invertible map and exact matching
to the source object, while exposing the cost of that map. They do not
prove a uniform weighted-class estimate for $L$ or for the direct
candidate $L=I$. The source's particular pulse retains its own source
estimates because of the equality \eqref{tgframe:exactpulsematch}.
Transport of the full viscosity-one witness to another physical
viscosity uses the already proved affine orthogonal and viscosity map.
Alternatively one may replace $\delta_1$ by $\delta_\nu$ in this
finite ODE construction, but the resulting changed-background
homogeneous pulse is not asserted to be the source's published witness.
